%% ----------------------------------------------------------------------------
%%  §1.9  Le funzioni trigonometriche
%% ----------------------------------------------------------------------------
\section{Le funzioni trigonometriche}
\label{sec:trigonometriche}

Riportiamo in questo paragrafo un breve riassunto di nozioni di trigonometria, in
genere già note al lettore; lo scopo è fissare le notazioni e, soprattutto, mettere in
evidenza le poche proprietà che useremo davvero nel seguito.

\subsection{La misura degli angoli in radianti}

Il lettore sa misurare gli angoli in gradi: un angolo retto misura $90^{\circ}$, un
angolo piatto $180^{\circ}$, un angolo giro $360^{\circ}$. Per studiare le funzioni
trigonometriche è opportuno adottare una diversa unità di misura.

\begin{definizione}[Radiante]\index{radiante}
La \emph{misura in radianti} di un angolo piano è la \emph{lunghezza dell'arco di
circonferenza di raggio $1$ e centro nel vertice dell'angolo intercettato dalle due
semirette che individuano l'angolo} (\cref{fig:radianti}).
\end{definizione}

Si conviene di denotare con $\pi$ (\emph{pi greco}) la lunghezza di una
semicirconferenza di raggio $1$. Ciò significa che, espresso in radianti, un angolo
piatto misura $\pi$, un angolo retto misura $\pi/2$, mentre un angolo giro misura
$2\pi$ (la lunghezza di una circonferenza di raggio $1$ è $2\pi$). Come tutti sanno, un
valore numerico approssimato di $\pi$ è $3{,}14\ldots$; la conversione tra gradi e
radianti è una proporzione: $\alpha^{\circ}:360^{\circ}=x:2\pi$.

\begin{idea}[Perché i radianti]
La scelta dei radianti non è un vezzo dei matematici. Il grado è un'unità arbitraria
(360 è un'eredità dell'astronomia babilonese), mentre il radiante è una misura
\emph{intrinseca}: un rapporto tra due lunghezze (arco e raggio), dunque un numero puro.
La ricompensa arriverà nel calcolo differenziale: la derivata di $\sen x$ è $\cos x$
\emph{solo se} $x$ è in radianti; in gradi comparirebbe un fastidioso fattore
$\pi/180$. Per un ingegnere il radiante è anche l'unità naturale per le velocità
angolari ($\omega$ in rad/s) e per la relazione $v=\omega r$ tra velocità periferica e
raggio.
\end{idea}

Analogamente a quanto si fa per l'ascissa su una retta, si definisce un'origine e un
verso di rotazione positivo (si suole scegliere il verso antiorario a partire dall'asse
delle $x$) e si considerano in modo naturale anche angoli maggiori di $2\pi$ radianti, o
angoli minori di zero. Così l'angolo $x$ geometricamente corrisponde all'angolo
$x+2\pi$ e anche all'angolo $x+4\pi$, oppure all'angolo $x-2\pi$, o in generale
all'angolo $x+2k\pi$, per ogni $k\in\Z$.

\begin{definizione}[Seno e coseno]\index{seno}\index{coseno}
Le funzioni $\sen x$ e $\cos x$ si definiscono rispettivamente come \emph{l'ordinata} e
\emph{l'ascissa} del punto che si trova sulla circonferenza di centro l'origine e raggio
$1$ e che sottende un angolo orientato di lunghezza $x$, a partire dall'asse delle
ascisse (\cref{fig:radianti}).
\end{definizione}

\begin{figure}[htb]
\centering
\begin{tikzpicture}[scale=1.55]
  \draw[-Latex, grigio] (-1.3,0) -- (1.5,0) node[below, font=\scriptsize] {$x$};
  \draw[-Latex, grigio] (0,-1.3) -- (0,1.5) node[left, font=\scriptsize] {$y$};
  \draw[grigio!70] (0,0) circle (1);
  \node[below right, font=\scriptsize, grigio] at (1,0) {$1$};
  \pgfmathsetmacro{\ang}{40}
  \draw[rame, very thick] (1,0) arc (0:\ang:1) node[midway, right=2pt, font=\scriptsize] {$x$};
  \draw[blu, thick] (0,0) -- (\ang:1);
  \fill[blu] (\ang:1) circle (1.2pt) node[above right, font=\scriptsize] {$P=(\cos x,\sen x)$};
  \draw[tratteggio] (\ang:1) -- ({cos(\ang)},0);
  \draw[tratteggio] (\ang:1) -- (0,{sin(\ang)});
  \draw[verde, very thick] (0,0) -- ({cos(\ang)},0) node[midway, below, font=\scriptsize] {$\cos x$};
  \draw[verde, very thick] ({cos(\ang)},0) -- ({cos(\ang)},{sin(\ang)}) node[pos=0.35, left, font=\scriptsize] {$\sen x$};
\end{tikzpicture}\hspace{8mm}
\begin{tikzpicture}[scale=1.55]
  \draw[-Latex, grigio] (-1.3,0) -- (1.5,0) node[below, font=\scriptsize] {$x$};
  \draw[-Latex, grigio] (0,-1.3) -- (0,1.5) node[left, font=\scriptsize] {$y$};
  \draw[grigio!70] (0,0) circle (1);
  \pgfmathsetmacro{\ang}{140}
  \draw[rame, very thick] (1,0) arc (0:\ang:1) node[pos=0.5, above right=1pt, font=\scriptsize] {$x$};
  \draw[blu, thick] (0,0) -- (\ang:1);
  \fill[blu] (\ang:1) circle (1.2pt) node[above left, font=\scriptsize] {$P$};
  \draw[tratteggio] (\ang:1) -- ({cos(\ang)},0);
  \draw[verde, very thick] (0,0) -- ({cos(\ang)},0) node[midway, below, font=\scriptsize] {$\cos x<0$};
  \draw[verde, very thick] ({cos(\ang)},0) -- ({cos(\ang)},{sin(\ang)}) node[midway, left, font=\scriptsize] {$\sen x$};
\end{tikzpicture}
\caption{Definizione di $\sen x$ e $\cos x$ sulla circonferenza di raggio $1$: $x$ è la
lunghezza dell'arco (in rame) percorso in verso antiorario a partire dal punto $(1,0)$.
A destra, un angolo del secondo quadrante: l'ascissa, e quindi il coseno, è negativa.}
\label{fig:radianti}
\end{figure}

Le funzioni $\sen x$, $\cos x$ sono definite per ogni $x\in\R$, mentre l'immagine delle
due funzioni è compresa tra $-1$ e $1$, cioè:
\begin{equation}\label{eq:sen-cos-limitate}
  -1\le\sen x\le1, \qquad -1\le\cos x\le1, \qquad\perogni x\in\R.
\end{equation}
Inoltre, poiché agli angoli $x$ e $x+2\pi$ corrisponde lo stesso punto $P$, le due
funzioni sono \emph{periodiche} di periodo $2\pi$:
$\sen(x+2\pi)=\sen x$ e $\cos(x+2\pi)=\cos x$. Il grafico delle due funzioni è riportato
nella \cref{fig:sen-cos}. Nel capitolo sullo studio di funzione sarà indicato come
ottenere il grafico; per ora il lettore controlli sul grafico il segno delle due funzioni
in base alla definizione. Le funzioni $\sen x$ e $\cos x$ non sono monotòne su $\R$.

\begin{figure}[htb]
\centering
\begin{tikzpicture}
\begin{axis}[width=\textwidth, height=4.8cm, xmin=-6.9, xmax=6.9, ymin=-1.4, ymax=1.5,
  xtick={-6.2832,-3.1416,-1.5708,1.5708,3.1416,4.7124,6.2832},
  xticklabels={$-2\pi$,$-\pi$,$-\frac{\pi}{2}$,$\frac{\pi}{2}$,$\pi$,$\frac{3}{2}\pi$,$2\pi$},
  ytick={-1,1}, samples=300, axis on top, xticklabel style={fill=white, inner sep=1pt}]
  \addplot[curva, domain=-6.8:6.8] {sin(deg(x))};
  \addplot[curva2, dashed, domain=-6.8:6.8] {cos(deg(x))};
  \node[blu, font=\scriptsize] at (axis cs:2.3,1.2) {$\sen x$};
  \node[rame, font=\scriptsize] at (axis cs:-1.0,1.2) {$\cos x$};
\end{axis}
\end{tikzpicture}
\caption{I grafici di $\sen x$ (linea continua) e $\cos x$ (tratteggiata). Il coseno è
il seno «in anticipo» di $\pi/2$: $\cos x=\sen(x+\pi/2)$.}
\label{fig:sen-cos}
\end{figure}

\subsection{La relazione fondamentale e il teorema di Pitagora}

Delle numerose relazioni tra queste due funzioni, la più importante è senza dubbio
quella che segue dal teorema di Pitagora sui triangoli rettangoli. È utile ricordare
come si ottiene una dimostrazione del teorema di Pitagora, perché è un piccolo
capolavoro di economia.

\begin{teorema}[Pitagora]\label{teo:pitagora}\index{teorema di Pitagora}
In un triangolo rettangolo con cateti lunghi $a$, $b$ e ipotenusa lunga $c$ risulta
\begin{equation}\label{eq:pitagora}
  a^{2}+b^{2}=c^{2}.
\end{equation}
\end{teorema}

\begin{schema}
\Ipotesi{Triangolo rettangolo con cateti lunghi $a$, $b$ e ipotenusa lunga $c$.}
\Tesi{$a^{2}+b^{2}=c^{2}$.}
\end{schema}

\begin{strategia}
L'idea è di calcolare \emph{una stessa area in due modi}. Disponiamo quattro copie del
triangolo dentro un quadrato di lato $a+b$ come nella \cref{fig:pitagora}: la regione
che resta scoperta è un quadrilatero di lato $c$, e un piccolo ragionamento sugli
angoli mostra che è un quadrato. Allora l'area del quadrato grande, $(a+b)^{2}$, è
uguale a quattro triangoli più il quadrato di lato $c$: sviluppando e semplificando
resta proprio la \eqref{eq:pitagora}.
\end{strategia}

\begin{dimostrazione}
Si veda la \cref{fig:pitagora}. Il triangolo rettangolo, ripetuto quattro volte, ha
cateti lunghi $a$, $b$ e angoli $\alpha$, $\beta$, oltre all'angolo retto. Evidentemente
$\alpha+\beta=\pi/2$ (perché la somma degli angoli interni del triangolo vale $\pi$);
ne segue che anche l'angolo $\gamma$ in figura è uguale a $\pi/2$, essendo
$\alpha+\gamma+\beta=\pi$ (angolo piatto). Perciò il rombo di lato $c$ rappresentato nella
parte destra della figura è in realtà un quadrato. L'area del quadrato più grande, di
lato $a+b$, vale $(a+b)^{2}$. Tale area si può ottenere sommando le aree delle figure
componenti (quattro triangoli rettangoli di cateti $a$, $b$, e un quadrato di lato
$c$), cioè $4\cdot(ab/2)+c^{2}$. Perciò
\[
  (a+b)^{2}=a^{2}+2ab+b^{2}=2ab+c^{2},
\]
da cui, semplificando $2ab$ (\cref{prop:semplificazione-somma}), il teorema di Pitagora
\eqref{eq:pitagora}.
\end{dimostrazione}

\begin{figure}[htb]
\centering
\begin{tikzpicture}[scale=0.85, every node/.style={font=\scriptsize}]
  \pgfmathsetmacro{\a}{1.4}\pgfmathsetmacro{\b}{2.4}
  \pgfmathsetmacro{\s}{\a+\b}
  % quadrato di lato a+b con il quadrato inclinato di lato c
  \draw[grigio, thick] (0,0) rectangle (\s,\s);
  \fill[bluchiaro] (\a,0) -- (\s,\a) -- (\b,\s) -- (0,\b) -- cycle;
  \draw[blu, thick] (\a,0) -- (\s,\a) -- (\b,\s) -- (0,\b) -- cycle;
  \draw[rame, thick] (0,0) -- (\a,0) -- (0,\b) -- cycle;
  \node[below] at (\a/2,0) {$a$}; \node[below] at (\a+\b/2,0) {$b$};
  \node[left] at (0,\b/2) {$b$}; \node[left] at (0,\b+\a/2) {$a$};
  \node[right] at (\s,\a/2) {$a$}; \node[right] at (\s,\a+\b/2) {$b$};
  \node[above] at (\b/2,\s) {$b$}; \node[above] at (\b+\a/2,\s) {$a$};
  \node[blu] at (\s/2,\s/2) {$c^{2}$};
  % angoli
  \draw[rame] (\a,0) ++(180:0.35) arc (180:120:0.35); \node[rame] at (\a-0.5,0.25) {$\alpha$};
  \draw[rame] (0,\b) ++(270:0.35) arc (270:300:0.35); \node[rame] at (0.15,\b-0.6) {$\beta$};
  \draw[blu] (\a,0) ++(0:0.35) arc (0:30.26:0.35); \node[blu] at (\a+0.55,0.2) {$\beta$};
  \node[blu] at (\a,0.45) {$\gamma$};
\end{tikzpicture}
\caption{La dimostrazione «per aree» del teorema di Pitagora. Il quadrato grande ha lato
$a+b$; i quattro triangoli rettangoli (uno è evidenziato in rame) lasciano scoperto il
quadrato azzurro di lato $c$. Poiché $\alpha+\beta=\pi/2$, l'angolo $\gamma$ è retto.}
\label{fig:pitagora}
\end{figure}

Applicando il teorema di Pitagora al triangolo rettangolo di cateti $\abs{\sen x}$ e
$\abs{\cos x}$ e ipotenusa uguale a $1$ (\cref{fig:radianti}), si trova la
\emph{relazione fondamentale} (se il triangolo degenera, cioè se $\sen x=0$ oppure
$\cos x=0$, il punto $P$ sta su un asse, l'altra coordinata vale $\pm1$ e la relazione
è ovvia)\index{relazione fondamentale della trigonometria}
\begin{equation}\label{eq:relazione-fondamentale}
  \sen^{2}x+\cos^{2}x=1, \qquad\perogni x\in\R.
\end{equation}
Sono anche importanti le \emph{formule di addizione}\index{formule di addizione}, che
riportiamo senza dimostrazione:
\begin{align}
  \sen(x_{1}\pm x_{2}) &= \sen x_{1}\cos x_{2}\pm\sen x_{2}\cos x_{1}, \label{eq:addizione-sen}\\
  \cos(x_{1}\pm x_{2}) &= \cos x_{1}\cos x_{2}\mp\sen x_{1}\sen x_{2}. \label{eq:addizione-cos}
\end{align}
Se nelle identità sopra scritte scegliamo il segno $+$ e poniamo $x_{1}=x_{2}=x$,
otteniamo le \emph{formule di duplicazione}:
\begin{equation}\label{eq:duplicazione}
  \sen2x=2\sen x\cos x, \qquad \cos2x=\cos^{2}x-\sen^{2}x .
\end{equation}
Dalle \eqref{eq:addizione-sen} e \eqref{eq:addizione-cos} con $x_{1}=0$ e il segno $-$
si ricava anche che il coseno è una funzione \emph{pari} e il seno è \emph{dispari}: $\cos(-x)=\cos x$,
$\sen(-x)=-\sen x$, come si vede del resto dalla simmetria della circonferenza rispetto
all'asse $x$.

\subsection{La funzione tangente}

\begin{definizione}[Tangente]\index{tangente}
A partire dalle funzioni $\sen x$ e $\cos x$ si definisce la \emph{funzione tangente}
\begin{equation}\label{eq:tangente}
  \tg x=\frac{\sen x}{\cos x},
\end{equation}
definita se $\cos x\neq0$, cioè se $x\neq\pi/2+k\pi$, con $k\in\Z$.
\end{definizione}

\begin{figure}[htb]
\centering
\begin{tikzpicture}
\begin{axis}[width=7.2cm, height=5.6cm, xmin=-5, xmax=5, ymin=-3.5, ymax=3.5,
  xtick={-4.7124,-3.1416,-1.5708,1.5708,3.1416,4.7124},
  xticklabels={$-\frac32\pi$,$-\pi$,$-\frac{\pi}{2}$,$\frac{\pi}{2}$,$\pi$,$\frac32\pi$},
  ytick=\empty, samples=200, axis on top, xticklabel style={fill=white, inner sep=1pt}]
  \foreach \k in {-1,0,1} {
    \addplot[curva, domain={-1.29+\k*3.1416}:{1.29+\k*3.1416}] {tan(deg(x))};}
  \foreach \v in {-4.7124,-1.5708,1.5708,4.7124} {\addplot[tratteggio] coordinates {(\v,-3.5) (\v,3.5)};}
  \node[font=\scriptsize, blu] at (axis cs:0.6,2.6) {$\tg x$};
\end{axis}
\end{tikzpicture}\hfill
\begin{tikzpicture}[scale=1.5]
  \draw[-Latex, grigio] (-0.3,0) -- (1.9,0);
  \draw[-Latex, grigio] (0,-0.3) -- (0,1.6);
  \draw[grigio!70] (1,0) arc (0:90:1);
  \pgfmathsetmacro{\ang}{50}
  \draw[blu, thick] (0,0) -- ({1/cos(\ang)*cos(\ang)},{1/cos(\ang)*sin(\ang)});
  \draw[rame, very thick] (1,0) arc (0:\ang:1) node[midway, right=2pt, font=\scriptsize] {$x$};
  \fill[blu] (\ang:1) circle (1.1pt) node[above left, font=\scriptsize] {$P$};
  \draw[tratteggio] (\ang:1) -- ({cos(\ang)},0) node[below, font=\scriptsize] {$A$};
  \draw[verde, very thick] (1,0) -- (1,{tan(\ang)}) node[midway, right, font=\scriptsize] {$\tg x$};
  \fill[verde] (1,{tan(\ang)}) circle (1.1pt) node[above right, font=\scriptsize] {$T$};
  \node[below, font=\scriptsize] at (1,0) {$B$};
  \node[below left, font=\scriptsize] at (0,0) {$O$};
  \node[left, font=\scriptsize] at (0,1) {$1$};
  \draw[grigio] (0,1) -- (-0.04,1);
  \node[font=\scriptsize, grigio, anchor=east] at ({cos(\ang)-0.02},{sin(\ang)*0.2}) {$\sen x$};
\end{tikzpicture}
\caption{A sinistra: il grafico di $\tg x$, con gli asintoti verticali
$x=\pi/2+k\pi$. A destra: la tangente come lunghezza del segmento $BT$ sulla retta
tangente alla circonferenza nel punto $B=(1,0)$.}
\label{fig:tangente}
\end{figure}

Usando le proprietà dei triangoli simili ($OAP$ e $OBT$ nella \cref{fig:tangente}),
$BT=BT/OB=AP/OA=\sen x/\cos x$, si vede che la tangente di un angolo $x$ si può
rappresentare come lunghezza del segmento $BT$. Dalla figura risulta anche chiaro che,
se $0<x<\pi/2$, allora valgono le disuguaglianze
\begin{equation}\label{eq:sen-x-tg}
  0<\sen x<x<\tg x,
\end{equation}
perché la corda è più corta dell'arco (il segmento $AP$ è un cateto del triangolo
rettangolo $APB$, quindi è più corto dell'ipotenusa $PB$, cioè della corda, a sua volta
più corta dell'arco $x$) e l'area del settore circolare $OBP$, che vale $x/2$, è
minore dell'area del triangolo $OBT$, che vale $\tg x/2$. La \eqref{eq:sen-x-tg} è
l'unico fatto «geometrico» di cui avremo bisogno nel seguito: da essa discende il
limite fondamentale $\lim_{x\to0}(\sen x)/x=1$, e da questo tutte le derivate delle
funzioni trigonometriche. La tangente è periodica di periodo $\pi$ e, in ciascun
intervallo $(-\pi/2+k\pi,\pi/2+k\pi)$, è strettamente crescente.

Riportiamo infine una tabella con i valori delle funzioni trigonometriche per alcuni
angoli di uso frequente. Il modo con cui ottenere tali valori è indicato
nell'\cref{es:valori-notevoli}.

\begin{table}[htb]
\centering
\renewcommand{\arraystretch}{1.35}
\setlength{\tabcolsep}{5pt}
\begin{tabular}{l*{8}{c}}
\toprule
$x$ (radianti) & $0$ & $\pi/6$ & $\pi/4$ & $\pi/3$ & $\pi/2$ & $\pi$ & $\tfrac32\pi$ & $2\pi$ \\
$x$ (gradi) & $0$ & $30$ & $45$ & $60$ & $90$ & $180$ & $270$ & $360$ \\
\midrule
$\sen x$ & $0$ & $\tfrac12$ & $\tfrac{\sqrt2}{2}$ & $\tfrac{\sqrt3}{2}$ & $1$ & $0$ & $-1$ & $0$ \\
$\cos x$ & $1$ & $\tfrac{\sqrt3}{2}$ & $\tfrac{\sqrt2}{2}$ & $\tfrac12$ & $0$ & $-1$ & $0$ & $1$ \\
$\tg x$ & $0$ & $\tfrac{\sqrt3}{3}$ & $1$ & $\sqrt3$ & non def. & $0$ & non def. & $0$ \\
\bottomrule
\end{tabular}
\caption{Valori notevoli delle funzioni trigonometriche.}
\label{tab:valori-notevoli}
\end{table}

\begin{esercizio}\label{es:valori-notevoli}
Ricavare i valori della \cref{tab:valori-notevoli} per $x=\pi/4$ e $x=\pi/6$.
\begin{soluzione}
Per $x=\pi/4$ il triangolo $OAP$ della \cref{fig:radianti} è rettangolo isoscele,
quindi $\sen x=\cos x$; dalla relazione fondamentale $2\sen^{2}x=1$, cioè
$\sen x=\cos x=\sqrt2/2$ (positivi perché $P$ è nel primo quadrante), e $\tg x=1$. Per
$x=\pi/6$: il triangolo di vertici $O$, $P$ e il simmetrico $P'$ di $P$ rispetto
all'asse $x$ ha due lati uguali a $1$ e l'angolo in $O$ uguale a $\pi/3$: è
equilatero, dunque $PP'=1$ e $\sen(\pi/6)=PP'/2=1/2$. Dalla relazione fondamentale
$\cos(\pi/6)=\sqrt{1-1/4}=\sqrt3/2$ e $\tg(\pi/6)=(1/2)/(\sqrt3/2)=1/\sqrt3=\sqrt3/3$.
I valori per $\pi/3$ seguono scambiando seno e coseno, perché $\pi/3=\pi/2-\pi/6$ e
$\sen(\pi/2-x)=\cos x$ (formula di addizione). \emph{Scelta logica:} i valori notevoli
si ricavano da figure con molte simmetrie (il triangolo rettangolo isoscele, il triangolo
equilatero), nelle quali una relazione geometrica elementare fornisce un'equazione per
$\sen x$; la relazione fondamentale dà poi il coseno, con il segno fissato dal quadrante.
\end{soluzione}
\end{esercizio}

\begin{python}[Verifica delle identità trigonometriche e delle disuguaglianze \eqref{eq:sen-x-tg}]
import math
# Relazione fondamentale e duplicazione, per angoli "a caso"
for x in (0.3, 1.0, 2.5, -4.0):
    print(math.isclose(math.sin(x)**2 + math.cos(x)**2, 1.0),
          math.isclose(math.sin(2*x), 2*math.sin(x)*math.cos(x)))

# Le disuguaglianze 0 < sen x < x < tg x per 0 < x < pi/2
for x in (0.1, 0.5, 1.0, 1.5):
    s, t = math.sin(x), math.tan(x)
    print(f"x={x:.1f}:  sen x={s:.4f} < x={x:.4f} < tg x={t:.4f}")

# I valori notevoli: math.radians converte i gradi in radianti
for gradi in (0, 30, 45, 60, 90):
    x = math.radians(gradi)
    print(f"{gradi:3d}°  sen={math.sin(x):.4f}"
          f"  cos={math.cos(x):.4f}")
\end{python}
\begin{risultato}
True True
True True
True True
True True
x=0.1:  sen x=0.0998 < x=0.1000 < tg x=0.1003
x=0.5:  sen x=0.4794 < x=0.5000 < tg x=0.5463
x=1.0:  sen x=0.8415 < x=1.0000 < tg x=1.5574
x=1.5:  sen x=0.9975 < x=1.5000 < tg x=14.1014
  0°  sen=0.0000  cos=1.0000
 30°  sen=0.5000  cos=0.8660
 45°  sen=0.7071  cos=0.7071
 60°  sen=0.8660  cos=0.5000
 90°  sen=1.0000  cos=0.0000
\end{risultato}

\begin{attenzione}[$\cos(90^{\circ})$ non è esattamente $0$ al calcolatore]
Nell'ultima riga dell'output il coseno di $90^{\circ}$ appare come $0{,}0000$, ma il
valore calcolato da Python è $6{,}1\cdot10^{-17}$: il numero $\pi/2$ non è
rappresentabile esattamente in virgola mobile, e il calcolatore sta valutando il coseno
di un numero \emph{vicinissimo} a $\pi/2$, ma non uguale. Confrontare numeri in virgola
mobile con \texttt{==} è quasi sempre un errore; per questo nel codice usiamo
\texttt{math.isclose}.
\end{attenzione}
